\section{Support-Sumset Geometry}

\paragraph{Digit support.}
For a digit vector \(a=(a_i)\), its support is
\[
    S_X = \{i : a_i \ne 0\}.
\]
Similarly \(S_Y=\{j:b_j\ne 0\}\).

\paragraph{Support sumset.}
The support sumset is
\[
    S_X + S_Y = \{i+j : i\in S_X,\ j\in S_Y\}.
\]
It is the degree-level footprint of all possible digit interactions before
carry normalization.

\paragraph{Representation counts.}
For each degree \(k\), define
\[
    \rho_{X,Y}(k) = \#\{(i,j)\in S_X\times S_Y : i+j=k\}.
\]
The number \(\rho_{X,Y}(k)\) is the diagonal density of the support geometry. The
additive energy is
\[
    E(S_X,S_Y)=\sum_k \rho_{X,Y}(k)^2,
\]
which measures concentration of representation counts across anti-diagonals.

\begin{proposition}[Raw support as a sumset]
Assume the digits are nonnegative. Let
\[
    c_k=\sum_{i+j=k} a_i b_j.
\]
Then
\[
    \operatorname{supp}(c)=S_X+S_Y.
\]
\end{proposition}

\begin{proof}
If \(k\in S_X+S_Y\), then there exist \(i\in S_X\) and \(j\in S_Y\) with
\(i+j=k\). Since digits are nonnegative and supported digits are positive,
\(a_i b_j>0\), so \(c_k>0\). Conversely, if \(c_k>0\), at least one summand
\(a_i b_j\) with \(i+j=k\) is positive. Hence \(i\in S_X\), \(j\in S_Y\),
and \(k\in S_X+S_Y\).
\end{proof}

\begin{proposition}[Representation-count bound]
For base \(B\), every raw coefficient satisfies
\[
    c_k \le (B-1)^2 \rho_{X,Y}(k).
\]
\end{proposition}

\begin{proof}
Every digit product satisfies \(a_i b_j\le (B-1)^2\). At degree \(k\), there
are exactly \(\rho_k\) support pairs that can contribute nonzero products.
Summing these products gives the stated bound.
\end{proof}

Carry complexity is influenced by concentration of \(\rho_{X,Y}(k)\). Large
representation counts and high additive energy indicate many support-pair
collisions on the same anti-diagonals, increasing the potential raw mass that
carry normalization must absorb. The realized carry still depends on digit
values and incoming carry, but support-sumset geometry provides a structural
predictor for where nonlinear normalization may become active.
